\documentclass[12pt,a4paper]{article}
\input{../dipecs-report-preamble}

\renewcommand{\reporttitle}{DiPECS 可行性研究报告}
\renewcommand{\reportsubtitle}{端侧 AIOS 的数据链路、治理边界与实现路径}
\renewcommand{\reporttype}{可行性研究报告}
\renewcommand{\reportdate}{2026 年 4 月 15 日}

\begin{document}

\pagenumbering{gobble}
\makecover
\clearpage
\pagenumbering{arabic}

{\hypersetup{linkcolor=black}\tableofcontents}
\newpage

\begin{dipecsabstract}
\noindent\textbf{\zihao{-4} 摘要：}\zihao{-5}
\textbf{目的：}评估 DiPECS 端侧 AIOS 原型在 Android/Linux 环境下的技术可行性，明确建立可验证数据链路、隐私边界与动作治理边界的可行条件。
\textbf{方法：}从数据采集、系统架构、决策模型、动作执行、隐私约束和实验评估六个方面，结合 Android 公开 API、Rust 核心实现、离线 replay 和项目代码仓库进行可行性论证。
\textbf{结果：}Android 公开 API 与 Rust 核心链路足以支撑一条可脱敏、可聚合、可审查、可回放的端侧 AIOS 原型；高权限采集、真实系统预加载和复杂语义模型仍是后续阶段需要逐步推进的难点。
\textbf{结论：}DiPECS 可以先以本地规则基线、公开 API 数据源和受控动作边界实现可运行原型，再逐步接入云端模型、复杂预测后端和高权限采集路径。

\vspace{0.5em}
\noindent\textbf{\zihao{-4} 关键词：}\zihao{-5} 可行性分析；端侧智能系统；隐私空气隙；动作治理；Android 公开 API
\end{dipecsabstract}

\section{引言}\label{sec:intro}

DiPECS 的目标不是实现一个通用聊天助手，而是在 Android/Linux 场景下验证端侧 AIOS 的核心机制：
系统如何从多源事件中获得上下文，如何在本地完成隐私脱敏，如何把上下文交给本地规则或云端模型进行意图判断，以及如何确保模型建议不会绕过本地策略直接执行。
AIOS 相关工作已经说明，LLM agent 需要调度、上下文、存储和访问控制等系统级支撑；DiPECS 则把这个问题收敛到 Android/Linux 端侧场景 \cite{aios}。

当前项目采用 Kotlin + Rust 的工程路线。
Kotlin 负责 Android 应用层能力，包括权限申请、通知监听、前台服务、设备状态采集和动作桥接；Rust 负责统一事件模型、隐私空气隙、窗口聚合、决策路由、策略审查、动作生命周期、离线 replay 和审计记录。
两侧通过 JSONL/schema 边界传输结构化数据，使同一套 Rust 核心逻辑既能处理真实 Android trace，也能处理离线测试夹具。

系统主链路可以概括为：
\texttt{RawEvent} $\rightarrow$ \texttt{PrivacyAirGap} $\rightarrow$ \texttt{SanitizedEvent} $\rightarrow$
\texttt{StructuredContext} $\rightarrow$ \texttt{DecisionRouter} $\rightarrow$ \texttt{PolicyEngine} $\rightarrow$
\texttt{ActionLifecycle} $\rightarrow$ \texttt{AuditRecord}。
这条链路中的每一步都有明确输入输出，因此项目可以通过分阶段开发方式逐步实现：先跑通离线 replay 和本地规则，再接入 Android 公共 API trace，最后补充云端模型、Android bridge 和更复杂的预测后端。

\section{系统设计可行性}\label{sec:architecture}

\subsection{总体架构}\label{subsec:architecture-overview}

DiPECS 的架构可以拆为六层，表 \ref{tab:layers} 给出了各层的代表模块与职责。

\begin{dipecstable}{DiPECS 分层设计}{tab:layers}
  \begin{tabular}{P{2.2cm}P{3.6cm}P{7.0cm}}
    \toprule
    \rowcolor{gray!15}\textbf{层级} & \textbf{代表模块} & \textbf{主要职责} \\
    \midrule
    采集层 & Android collector, aios-collector & 采集应用切换、通知、屏幕、电量、进程资源和 JSONL trace。 \\
    \dtr
    隐私层 & PrivacyAirGap & 将 raw event 转为 sanitized event，丢弃通知正文、文件路径等敏感信息。 \\
    \dtr
    上下文层 & WindowAggregator & 按时间窗口构造 StructuredContext 和 ContextSummary。 \\
    \dtr
    决策层 & DecisionRouter, RuleBased, CloudLlm & 根据复杂度、隐私分数和后端状态选择本地或云端决策。 \\
    \dtr
    治理层 & PolicyEngine, ActionLifecycle & 检查 capability、风险、置信度、target 和 action schema。 \\
    \dtr
    执行/验证层 & aios-action, OfflineAdapter, replay & 执行低风险动作或离线模拟，生成 audit record 和 hash。 \\
    \bottomrule
  \end{tabular}
\end{dipecstable}

这种分层的可行性来自两个事实。
第一，Android 和 Linux 已经能提供足够的低粒度事件，支持最小原型；
第二，Rust 中的核心模块不依赖 Android UI，可以用离线 trace 进行确定性测试。
即使 Android 物理设备能力尚未完全接入，系统仍可通过 replay 验证隐私、路由、策略和审计链路。

\subsection{数据采集可行性}\label{subsec:data-collection}

Android 普通应用无法任意读取其他应用内部状态，但仍可在用户授权下使用公开 API 获取有限而稳定的上下文。
UsageStats 可以提供应用前后台切换趋势；NotificationListener 可以获取通知来源、类别和出现时间；BatteryManager、ConnectivityManager 和 PowerManager 可以提供电量、网络和屏幕状态 \cite{android-usagestats,android-notification,android-battery}。
这些信号虽然不能完整表达用户意图，但足以构造第一阶段的行为窗口。

更高权限的数据源，如 Accessibility、ADB/Shizuku、Binder trace、eBPF 或 root 环境，可以提高观测精度，但部署成本和隐私风险显著更高。
因此本项目的可行路径是：默认主链路依赖公开 API 和本机 Linux 采集；高权限数据源只作为实验增强，不作为最小闭环的前提 \cite{android-accessibility,aosp-ebpf}。

\subsection{隐私脱敏可行性}\label{subsec:privacy}

端侧上下文天然包含敏感信息。
通知标题和正文可能包含联系人、验证码或工作内容；文件路径可能包含项目名、客户名或个人目录；精确地点可能直接暴露用户轨迹。
因此系统不能把 raw event 直接交给云端模型。

DiPECS 采用本地确定性脱敏。
\texttt{RawEvent} 只存在于 collector 到 core 的短路径中，经过 \texttt{PrivacyAirGap} 后，后续模块只处理 \texttt{SanitizedEvent}。
例如，通知正文被转换为文本长度、书写系统和 semantic hint；文件路径只用于扩展名类别判断；notification key 等高风险字段被丢弃。
这种设计牺牲一部分语义完整性，但换来清晰的隐私边界和可测试的泄漏防护。

\subsection{上下文聚合可行性}\label{subsec:aggregation}

单个事件很难支撑意图判断。
用户是否即将打开某个应用，通常要结合通知、前台应用、屏幕状态、电量和最近进程状态共同判断。
因此项目使用时间窗口聚合，把 \texttt{SanitizedEvent} 序列转为 \texttt{StructuredContext}。

窗口聚合的计算成本低，适合端侧运行；输出结构稳定，适合 replay；summary 字段又能降低后端决策复杂度。
这使得 RuleBased、CloudLlm 和未来的统计预测模型都可以共享同一输入格式，避免为不同后端重复处理 raw event。

\section{决策模型可行性}\label{sec:decision}

\subsection{RuleBased 基线}\label{subsec:rulebased}

RuleBased 是项目第一阶段后端的首选方案。
它不依赖网络、不需要训练数据，行为可解释，适合作为默认降级路径。
当前合理边界是：RuleBased 只生成低风险、能力范围内的动作，例如 \texttt{NoOp}、\texttt{KeepAlive} 和 \texttt{ReleaseMemory}。
它不应直接生成 \texttt{PreWarmProcess} 或 \texttt{PrefetchFile}，因为这些动作需要更明确的 target 语义和更高能力等级。

RuleBased 的价值不在于预测能力最强，而在于建立系统基线：给定同一 \texttt{StructuredContext}，它应产生稳定的 \texttt{IntentBatch}，并且不应持续触发 capability denial。
这个基线为后续云端模型和统计模型提供对照。

\subsection{统计预测模型}\label{subsec:statistical}

在 RuleBased 之后，最自然的增强是统计转移模型。
移动行为预测研究表明，用户应用切换存在明显时序规律。
项目可以从 replay trace 或真实 Android trace 中统计 foreground app 之间的转移概率，形成
$P(a_{t+1} \mid a_t)$ 或 $P(a_{t+1} \mid a_{t-k}, \ldots, a_t)$。

统计模型实现成本低，推理开销小，且容易解释。
它可以作为 LocalEvaluator 的第一版：在上下文中输出 top-k 候选 app 和置信度，再由 \texttt{PolicyEngine} 判断是否允许 \texttt{KeepAlive} 或更高级动作。
若统计模型不能显著优于 RuleBased，复杂序列模型也不应过早投入。

\subsection{云端 LLM 后端}\label{subsec:cloudllm}

云端 LLM 的可行性体现在语义理解和规则编排上。
对于包含多个 semantic hint 的复杂上下文，LLM 可以基于结构化 JSON 生成 \texttt{IntentBatch}，并给出风险等级、置信度和 rationale tags。
当前设计中，CloudLlm 是可选后端：只有配置完整且 router 判断上下文复杂度足够时才调用；失败后回退到 RuleBased，并由 circuit breaker 防止连续错误放大。
OS-Copilot/FRIDAY 这类系统表明，LLM 可以组织跨应用任务和工具调用；但在 DiPECS 中，这类能力必须被本地 policy 边界约束 \cite{os-copilot,os-copilot-page}。

这种设计降低了云端依赖风险。
网络不可用、API 失败或模型输出格式错误都不会阻塞本地链路。
但同时也说明，云端模型不能作为唯一控制面；所有云端输出仍必须经过本地策略和动作生命周期。

\section{动作执行可行性}\label{sec:action}

动作执行是项目风险最高的部分。
相比预测``用户可能打开某个应用''，执行``预热、保活、预取、释放内存''等动作会直接影响系统资源。
为了降低风险，DiPECS 将动作分为建议、授权和执行三个阶段。
表 \ref{tab:actions} 总结了各动作类型的可行性与限制。

\begin{dipecstable}{动作类型的可行性与限制}{tab:actions}
  \begin{tabular}{P{2.4cm}P{5.0cm}P{5.0cm}}
    \toprule
    \rowcolor{gray!15}\textbf{动作} & \textbf{当前可行语义} & \textbf{主要限制} \\
    \midrule
    NoOp & 不做动作，只记录系统观察结果 & 无直接收益，但适合安全降级。 \\
    \dtr
    KeepAlive & 对已出现 app 或系统任务做轻量保活提示 & Android 普通应用能做的真实保活能力有限。 \\
    \dtr
    ReleaseMemory & 释放系统或缓存类资源，降低压力 & package 级释放需与 Android bridge 能力对齐。 \\
    \dtr
    PreWarmProcess & 预热应用或本系统资源 & 第三方 app 预热受平台限制，不能默认承诺。 \\
    \dtr
    PrefetchFile & 对 url/uri 或可访问内容做预取 & 需要明确 target 权限和 policy 放行规则。 \\
    \bottomrule
  \end{tabular}
\end{dipecstable}

因此，最可行的执行路线是先用 OfflineAdapter 验证 action lifecycle，再逐步接入 Android bridge。
真实执行层应优先支持 Android-safe 的动作，例如本应用资源预热、公开 content URI 预取、cache trim 或用户可见通知提示。
对第三方进程强制预热、系统级内存调度等能力，应明确标记为 root/AOSP 阶段目标。

\section{验证与评估可行性}\label{sec:evaluation}

DiPECS 采用离线 replay 作为主要验证手段。
Replay 不访问真实 Android、不访问网络，也不执行真实系统动作，而是读取 JSONL trace，经过相同的 Rust pipeline，生成 intent、policy decision、execute record 和 audit hash。
这使项目可以在没有物理设备或云端的情况下持续验证核心逻辑。

移动端预加载系统的评估不能只看预测是否命中，还要看启动延迟、能耗和误预加载成本。
FALCON 对 app prelaunch 的评估方式可以作为后续实验设计参考 \cite{falcon}。
评估指标应分为三组：

\begin{enumerate}[label=(\arabic*)]
  \item \textbf{预测指标}：top-1/top-k 命中率、置信度校准、不同后端相对 RuleBased 的提升；
  \item \textbf{系统指标}：动作授权数、拒绝数、noop rate、policy denial reason、执行失败数、replay hash 稳定性；
  \item \textbf{资源指标}：动作触发后的启动延迟变化、CPU/内存开销、误预测导致的资源浪费率。
\end{enumerate}

第一阶段应优先使用系统指标和 replay 稳定性，因为它们直接验证系统边界是否正确。
只有在动作桥接更完整后，再重点评估启动延迟和资源收益。

\section{风险分析}\label{sec:risk}

\subsection{权限与部署风险}\label{subsec:risk-permission}

Android 公开 API 能力有限，高权限数据源部署复杂。
应对策略是把公开 API 作为默认主路径，把 root、eBPF、Binder trace 明确为实验增强。
报告、代码和展示都不应把高权限能力写成已经完成的默认依赖。

\subsection{隐私与语义损失风险}\label{subsec:risk-privacy}

脱敏会丢失通知正文、文件路径和应用内文本，可能降低模型判断能力。
应对策略是用 semantic hint、类别标签和 summary 字段补偿部分语义，同时用 replay/audit 确保脱敏边界不被绕过。

\subsection{云端不稳定风险}\label{subsec:risk-cloud}

云端 LLM 可能因为网络、配额或格式错误失败。
应对策略是保持 RuleBased fallback、配置 circuit breaker，并要求云端只输出结构化 JSON，不直接控制设备。

\subsection{动作收益不确定风险}\label{subsec:risk-benefit}

即使预测正确，动作也未必带来可感知收益。
例如某些应用启动本来很快，预加载反而浪费资源。
应对策略是引入 confidence threshold、预算限制和 wasted action rate，用实验数据决定哪些动作值得保留。

\section{阶段计划}\label{sec:plan}

表 \ref{tab:phases} 给出了分阶段开发路线。
表中日期为报告撰写时已知的计划节点，实际执行以代码仓库迭代为准。

\begin{dipecstable}{阶段性开发路线}{tab:phases}
  \begin{tabular}{P{1.8cm}P{5.0cm}P{5.8cm}}
    \toprule
    \rowcolor{gray!15}\textbf{阶段} & \textbf{目标} & \textbf{可验证结果} \\
    \midrule
    Phase 1 & 打通离线 replay 和 RuleBased 基线 & 同一 trace 产生稳定 intent、policy 和 audit hash。 \\
    \dtr
    Phase 2 & 接入 Android public-API trace & 模拟器或物理设备 JSONL 能进入 Rust pipeline。 \\
    \dtr
    Phase 3 & 完善 cloud/rule/local 后端路由 & 云端失败可回退，privacy-sensitive context 不上传。 \\
    \dtr
    Phase 4 & 增加统计预测后端 & 与 RuleBased 比较 top-k 命中率和 noop rate。 \\
    \dtr
    Phase 5 & 强化 Android-safe action bridge & 真实动作有审计记录，并能量化收益与开销。 \\
    \bottomrule
  \end{tabular}
\end{dipecstable}

\section{局限性与展望}\label{sec:limitations}

本可行性分析存在以下局限：

\begin{enumerate}[label=(\arabic*)]
  \item 分析以公开 API 能力和项目内部实现为主，尚未在真实物理设备上完成端到端长时间运行验证；
  \item 高权限数据源（Accessibility、Binder、eBPF）的可行性判断以能力研究为主，缺少稳定部署方案；
  \item 动作收益评估依赖后续实验，当前阶段无法给出启动延迟或能耗的量化结论；
  \item 云端 LLM 的可行性结论基于接口设计和假设场景，真实场景下的失败模式需进一步覆盖。
\end{enumerate}

后续工作应在守住数据链路、隐私边界和动作治理边界的前提下，逐步扩展：稳定 Android 到 Rust 的生产入口、建立 golden trace 回归基线、接入真实云端后端、验证 Android-safe 动作桥接，并在更长周期内评估预测模型和动作执行的真实收益。

\section{结论}\label{sec:conclusion}

综合分析，DiPECS 的最小闭环具备可行性。
Android 公开 API 足以提供第一阶段上下文；Rust 核心可以稳定承担脱敏、聚合、策略和 replay；RuleBased 可以作为本地安全基线；CloudLlm 和统计模型可以作为后续增强；ActionLifecycle 和 audit record 能限制动作风险并提供可验证记录。

项目当前需要避免过度承诺高权限采集和系统级资源调度。
可行的发展策略是先保证数据链路、隐私边界和动作治理正确，再逐步证明预测模型和真实动作是否带来性能收益。
这样既能形成可演示原型，也能保留继续扩展到 root/AOSP 环境的空间。

\bibliographystyle{unsrtnat}
\bibliography{../refs}

\end{document}
